\documentclass{article}
\usepackage{amsmath}
\usepackage{amssymb}
\usepackage{geometry}
\usepackage{graphicx}
\usepackage{xepersian}
\settextfont{XB Niloofar}

\geometry{a4paper, margin=1in}

\title{\lr{Feature Maps: Bridging to Kernel Methods} \\ (نگاشت‌های ویژگی: پلی به سمت روش‌های هسته)}
\author{}
\date{}
\begin{figure}[h!]
  \centering
  \includegraphics[width=0.75\textwidth]{img_page_2_1.png}
  \caption{نگاشت داده‌های غیرخطی در فضای اولیه به فضای ویژگی با بعد بالاتر برای تفکیک خطی}
\end{figure}

\begin{figure}[h!]
  \centering
  \includegraphics[width=0.75\textwidth]{img_page_3_1.png}
  \caption{حل مسئله خوشه‌بندی غیرخطی دو حلقه با نگاشت به فضای ویژگی با بعد بالاتر}
\end{figure}

\begin{figure}[h!]
  \centering
  \includegraphics[width=0.75\textwidth]{img_page_4_1.png}
  \caption{استفاده از نگاشت‌های ویژگی غیرخطی برای مسئله رگرسیون}
\end{figure}

\begin{figure}[h!]
  \centering
  \includegraphics[width=0.75\textwidth]{img_page_5_1.png}
  \caption{برازش تابع چندجمله‌ای غیرخطی بر روی داده‌های رگرسیون شبیه‌سازی‌شده}
\end{figure}

\begin{figure}[h!]
  \centering
  \includegraphics[width=0.75\textwidth]{img_page_6_1.png}
  \caption{نمایش عملکرد نامناسب رگرسیون خطی ساده بر روی داده‌های ذاتا غیرخطی}
\end{figure}

\begin{figure}[h!]
  \centering
  \includegraphics[width=0.75\textwidth]{img_page_8_1.png}
  \caption{پراکندگی داده‌ها پس از اعمال نگاشت ویژگی به فضای ویژگی دوبعدی غیرخطی}
\end{figure}














\begin{document}

\maketitle

\section*{فهرست مطالب}
\begin{enumerate}
    \item مقدمه‌ای بر نگاشت‌های ویژگی (\lr{Feature Maps})
    \item دسته‌بندی: مسئله دو حلقه
    \item خوشه‌بندی: مسئله دو حلقه
    \item مثال رگرسیون: استفاده از نگاشت‌های ویژگی برای داده‌های غیرخطی
\end{enumerate}

\newpage

\section*{۱. مقدمه‌ای بر نگاشت‌های ویژگی}
یک \lr{feature map} داده‌های ورودی را به فضایی با ابعاد بالاتر انتقال می‌دهد:
\[
\phi: \mathbb{R}^p \to \mathbb{R}^d \quad \text{که در آن} \quad d \gg p
\]

انگیزه کلیدی: قادر ساختن مدل‌های خطی برای حل مسائل غیرخطی از طریق:
\begin{itemize}
    \item نگاشت صریح برای دسته‌بندی/خوشه‌بندی
    \item انبساط پایه برای رگرسیون
\end{itemize}

\section*{۲. دسته‌بندی: مسئله دو حلقه}
\subsection*{فضای دوبعدی اصلی (شکست مدل خطی)}
نگاشت $\phi(x,y) = [x, y, x^2 + y^2]$ کلاس‌ها را با تبدیل فاصله شعاعی به یک ویژگی خطی، به صورت خطی تفکیک‌پذیر می‌کند.

\section*{۳. خوشه‌بندی: مسئله دو حلقه}
خوشه‌بندی با استفاده از نگاشت‌های ویژگی در مسئله دو حلقه انجام می‌شود تا داده‌ها در فضای منتقل‌شده تفکیک‌پذیر شوند.

\section*{۴. مثال رگرسیون: استفاده از نگاشت‌های ویژگی برای داده‌های غیرخطی}
\textbf{هدف:} نشان دادن اینکه چگونه نگاشت ویژگی‌ها به یک فضای با ابعاد بالاتر می‌تواند به یک مدل خطی اجازه دهد تا روی داده‌های غیرخطی برازش شود.

\textbf{انگیزه:} مدل‌های رگرسیون خطی استاندارد فرض می‌کنند که یک رابطه خطی بین ویژگی‌ها (متغیرهای مستقل) و هدف (متغیر وابسته) وجود دارد. وقتی رابطه زمینه‌ای غیرخطی باشد چه اتفاقی می‌افتد؟ یک مدل خطی ساده عملکرد ضعیفی خواهد داشت.

سناریویی را در نظر بگیرید که در آن داده‌ها از یک منحنی پیروی می‌کنند، به عنوان مثال، یک رابطه درجه دوم. یک خط مستقیم (حاصل از رگرسیون خطی) این منحنی را به طور موثر ثبت نمی‌کند.

\textbf{ایده راه حل:} ما می‌توانیم ویژگی‌های اصلی را به فضای با ابعاد بالاتر منتقل کنیم که در آن رابطه خطی می‌شود. این انتقال یک نگاشت ویژگی (\lr{feature map}) نامیده می‌شود که با $\Phi$ نشان داده می‌شود. بیایید این موضوع را با یک مثال نشان دهیم.

\subsection*{تولید داده‌های غیرخطی}
ما داده‌های مصنوعی ایجاد می‌کنیم که در آنها $y$ به طور درجه دوم به $x$ وابسته است، به علاوه مقداری نویز تصادفی برای واقعی‌تر کردن آن. داده‌ها به وضوح از یک منحنی پیروی می‌کنند، نه یک خط مستقیم.

\subsection*{تلاش برای رگرسیون خطی ساده (فضای اصلی)}
بیایید ببینیم یک مدل رگرسیون خطی استاندارد چگونه روی این داده‌ها بدون هیچ‌گونه تبدیل ویژگی عمل می‌کند.

مدل خطی تمام تلاش خود را می‌کند تا یک خط مستقیم را روی داده‌های منحنی برازش دهد، اما به وضوح برازش ضعیفی است. امتیاز $R^2$ احتمالاً کم خواهد بود، که نشان می‌دهد مدل مقدار زیادی از واریانس داده‌ها را توضیح نمی‌دهد.
\[
\text{Linear Regression Score } (R^2): 0.4260
\]

\subsection*{اعمال یک نگاشت ویژگی (ویژگی‌های چندجمله‌ای)}
اکنون، بیایید یک نگاشت ویژگی اعمال کنیم. از آنجا که می‌دانیم رابطه زمینه‌ای درجه دوم است ($y \approx a x^2 + b x + c$)، یک نگاشت ویژگی مناسب $\Phi$ ویژگی تک ما یعنی $x$ را به دو ویژگی تبدیل می‌کند: $x$ و $x^2$. بنابراین، $\Phi(x) = [x, x^2]$.

توجه داشته باشید که داده‌های ما اکنون در یک فضای ویژگی دو‌بعدی $[x, x^2]$ نمایش داده می‌شوند.

مدل خطی اصلی:
\[
\hat{y} = w^T x = \sum_{i=1}^{p} w_i x_i
\]

یک تبدیل $\phi$ که ویژگی‌ها را به ابعاد بالاتر نگاشت می‌کند:
\[
\phi: \mathbb{R}^p \to \mathbb{R}^d \quad (d \gg p)
\]

مدل جدید به صورت زیر در می‌آید:
\[
\hat{y} = w^T \phi(x)
\]
برای $p=1$ و درجه $=2$:
\[
[1, x_1] \to [1, x_1, x_1^2] \quad \text{توسط} \quad \phi
\]

\subsection*{رگرسیون خطی در فضای ویژگی با ابعاد بالاتر}
اکنون، ما یک مدل رگرسیون خطی را آموزش می‌دهیم، اما از ویژگی‌های تبدیل‌شده استفاده می‌کنیم. مدل وزن‌هایی را برای هر دو $x$ و $x^2$ یاد می‌گیرد، به طوری که به طور موثر مدلی به فرم زیر را برازش می‌دهد:
\[
y = w_1 x + w_2 x^2 + b
\]
که در آن $w_1$ و $w_2$ وزن‌ها (ضرایب) و $b$ عرض از مبدأ (بایاس) هستند. این یک مدل خطی نسبت به ویژگی‌های جدید $x$ و $x^2$ است.
\begin{itemize}
    \item امتیاز رگرسیون چندجمله‌ای: $\text{Polynomial Regression Score } (R^2): 0.8525$
    \item ضرایب ($w_1$, $w_2$): $[0.93366893, 0.56456263]$
    \item عرض از مبدأ ($b$): $[1.78134581]$
\end{itemize}

\subsection*{نتیجه‌گیری}
با نگاشت ویژگی تک اصلی $x$ به فضای با ابعاد بالاتر $[x, x^2]$، ما به یک مدل رگرسیون خطی استاندارد امکان دادیم تا رابطه غیرخطی (درجه دوم) موجود در داده‌ها را به طور کامل ثبت کند. برازش حاصل بسیار بهتر است، همانطور که به صورت بصری و با امتیاز به وضوح بالاتر $R^2$ نشان داده شده است.

این مثال قدرت نگاشت‌های ویژگی را نشان می‌دهد: تبدیل داده‌ها به فضایی که در آن مدل‌های خطی می‌توانند موثر واقع شوند، حتی اگر رابطه اصلی غیرخطی بوده باشد.

\textbf{انتقال به ترفند هسته (\lr{Transition to Kernel Trick}):} در حالی که این نگاشت ویژگی صریح برای داده‌های با ابعاد کم و تبدیل‌های ساده به خوبی کار می‌کند، محاسبه و ذخیره‌سازی این ویژگی‌های با ابعاد بالاتر می‌تواند از نظر محاسباتی پرهزینه یا حتی غیرممکن باشد، اگر داده‌های اصلی یا فضای ویژگی هدف ابعاد بسیار بالایی داشته باشند (مثلاً با استفاده از ویژگی‌های چندجمله‌ای با درجه بسیار بالا یا نگاشت‌های پیچیده دیگر). ترفند هسته (\lr{Kernel Trick}) یک میانبر ریاضی برای دستیابی به همان نتیجه کار در فضای ویژگی با ابعاد بالا بدون محاسبه صریح مختصات داده‌ها در آن فضا فراهم می‌کند. این روش به الگوریتم‌هایی که فقط به ضرب داخلی بین نقاط داده وابسته هستند (مانند ماشین‌های بردار پشتیبان - \lr{Support Vector Machines}) اجازه می‌دهد تا به طور ضمنی در فضای ویژگی با ابعاد بالا کار کنند و نگاشت‌های پیچیده را از نظر محاسباتی قابل مدیریت می‌سازد. این مثال رگرسیون به ایجاد این شهود کمک می‌کند که چرا چنین نگاشت‌های ضمنی مطلوب هستند.

\end{document}
